\section{No Global Black-Box Eliminator}
\label{sec:visions-philosophy-scoped-certification-principle}
\origin{ai}

\autoref{sec:visions-philosophy-black-box-decomposition} decomposes a
machine-learning black box into support candidates, certificate fields,
ledger gaps, debt, NDNA loci, optional verification hardening, and value
accounting. The present section records the corresponding limitation:
black-box resolution is scope-relative unless a global ledger is
displayed. A total feature label, total diagram, total proof, or total
interpretability narrative is not a global eliminator unless it carries
the behavior family, classifier family, NameCert family, global ledger,
verification ledger where applicable, and not-claimed boundary.

\begin{definition}[Explanation scope]
\label{def:philosophy-scoped-certification-explanation-scope}
An explanation scope is a tuple
$$
\Omega=(D,M_\theta,\mathcal A,B),
$$
where $D$ is the input or behavior source domain, $M_\theta$ is the
model state, $\mathcal A$ is the admitted perturbation, transformation,
transfer, or rewrite family, and $B$ is the behavior being explained. An
explanation is scoped when its certificate and ledger claims are made
only inside $\Omega$.
\end{definition}

\begin{definition}[Scoped black-box resolution]
\label{def:philosophy-scoped-black-box-resolution-principle}
Behavior $B$ is resolved on scope $\Omega$ when there is an explanation
object
$$
E_\Omega=(B,C,N_C,\Lambda_C)
$$
such that $C\Vdash_\Omega B$, $N_C=\NameCert(C)$, the ledger
$\Lambda_C$ is complete for $(B,C,\Omega)$, and the seal records
$\Omega$ together with the not-claimed boundary.
\end{definition}

\begin{definition}[Global behavior family]
\label{def:philosophy-global-behavior-family}
For a model state $M_\theta$, the global behavior family
$\mathcal B(M_\theta)$ is the family of behaviors being covered by a
global explanation claim. It may contain output behaviors, refusals,
errors, transfer behaviors, safety behaviors, intervention behaviors,
and boundary behaviors.
\end{definition}

\begin{definition}[Global resolution and global ledger]
\label{def:philosophy-global-resolution-ledger}
A model state is globally resolved when every
$B\in\mathcal B(M_\theta)$ has a closed explanation and these
explanations are covered by a common global ledger. A global ledger
$\Lambda_{\mathrm{global}}$ records, for every covered behavior and
supporting classifier, the required source, pattern, classifier,
stability, semantic, intervention, verification, generalization, and
not-claimed residue:
$$
\operatorname{Req}(B,C_B)
\subseteq
\operatorname{Rec}(\Lambda_{\mathrm{global}}).
$$
\end{definition}

\begin{definition}[Global black-box eliminator]
\label{def:philosophy-global-black-box-eliminator}
A global black-box eliminator is an explanation claim that asserts
global resolution of $M_\theta$. It is uncertified when it omits the
covered behavior family, the supporting classifier family, the NameCert
family, the global ledger, the verification ledger for verified rows, or
the not-claimed boundary.
\end{definition}

\begin{theorem}[Scoped resolution does not imply global resolution]
\label{thm:philosophy-scoped-resolution-not-global-resolution}
Resolving $B$ on a scope $\Omega$ does not imply global resolution of
$M_\theta$.
\end{theorem}
\begin{proof}
Scoped resolution supplies a closed explanation only for the target
$(B,\Omega)$. Global resolution requires closed explanations for every
behavior in $\mathcal B(M_\theta)$ and a common ledger covering their
required residue. The scoped certificate gives neither the remaining
behaviors nor their global ledger.
\end{proof}

\concretizedIn{ch:concrete-instances-scoped-explanation-ledger-namecert}{ScopedExplanationLedgerUp realizes scoped black-box certification as finite explanation-scope, classifier, ledger, refusal, transport, and NameCert rows.}

\begin{theorem}[Finite-sample closure is not open-domain closure]
\label{thm:philosophy-finite-sample-closure-not-open-domain-closure}
Closure on a finite sample domain $D_0$ does not by itself imply closure
on a larger open domain $D$.
\end{theorem}
\begin{proof}
The finite certificate records required residue for $D_0$. A larger
domain may contain boundary cases, out-of-distribution cases, semantic
ambiguity, and generalization residue not recorded by the finite ledger.
Thus the ledger inclusion for $D_0$ does not entail the ledger inclusion
for $D$.
\end{proof}

\begin{theorem}[Global resolution requires global ledger]
\label{thm:philosophy-global-resolution-requires-global-ledger-principle}
Any global resolution claim for $M_\theta$ requires a global ledger.
\end{theorem}
\begin{proof}
Global resolution requires each behavior in $\mathcal B(M_\theta)$ to
have a closed explanation. Each closed explanation requires ledger
coverage for its required residue. Taking the family of all such
required rows gives the required residue of the global claim, and a
ledger covering that family is precisely the global ledger.
\end{proof}

\begin{theorem}[No uncertified global black-box eliminator]
\label{thm:philosophy-no-uncertified-global-black-box-eliminator}
An uncertified global black-box eliminator is not a valid global
resolution claim.
\end{theorem}
\begin{proof}
By
\autoref{thm:philosophy-global-resolution-requires-global-ledger-principle},
global resolution requires a global ledger. It also requires behavior
coverage, supporting classifiers, NameCert rows, and the relevant
verification and not-claimed ledgers. An eliminator claim omitting these
rows lacks the conditions it asserts.
\end{proof}

\subsection{Classifier and Semantic Globalization}
\label{subsec:philosophy-classifier-semantic-globalization}

\begin{definition}[Single-classifier explanation]
\label{def:philosophy-single-classifier-explanation}
A single-classifier explanation uses one classifier $C^\star$ to explain
several behaviors. Two behaviors are classifier-separated when their
closed explanations require nonequivalent classifier relations unless a
product classifier or bridge ledger is supplied.
\end{definition}

\begin{definition}[Product classifier and bridge ledger]
\label{def:philosophy-product-classifier-bridge-ledger}
For classifiers $C_1,\ldots,C_n$, a product classifier is the joint
relation
$$
C_\Pi(x,y)
\Longleftrightarrow
\bigwedge_{i=1}^{n} C_i(x,y).
$$
A bridge ledger for such a product records each source scope, each
classifier boundary, the distinctions created by the product, the
behaviors unified or separated by it, and its not-claimed boundary.
\end{definition}

\begin{theorem}[A single classifier does not automatically explain many behaviors]
\label{thm:philosophy-single-classifier-not-many-behaviors}
If $B_1$ and $B_2$ are supported by nonequivalent classifiers, a closed
explanation of $B_1$ by $C_1$ does not automatically close an explanation
of $B_2$.
\end{theorem}
\begin{proof}
The certificate for $C_1\Vdash B_1$ records the source, classifier,
stability, and ledger rows for $B_1$. Nonequivalence of the classifiers
means that the public classification surface needed for $B_2$ can differ.
Without a proof that $C_1$ supports $B_2$, or a certified product or
bridge, the rows for $B_2$ are absent.
\end{proof}

\begin{theorem}[Product classifiers require bridge ledgers]
\label{thm:philosophy-product-classifiers-require-bridge-ledgers}
A product classifier used to unify several behavior explanations requires
a bridge ledger.
\end{theorem}
\begin{proof}
The product classifier changes classification granularity. It can split
old cells, expose boundary cases, and change which behavior rows are
covered by which classifier. These residues are created by the
unification operation itself and must be recorded for the unified claim
to close.
\end{proof}

\begin{definition}[Semantic globalizer]
\label{def:philosophy-semantic-globalizer}
A semantic globalizer is an explanation claim that uses one natural
language concept $K$ to unify several features, classifiers, behaviors,
or model states. It overreaches when it lacks a semantic ledger for
positive cases, negative cases, ambiguous cases, polysemy, domain
dependence, context dependence, label drift, unsupported connotations,
and not-claimed semantics.
\end{definition}

\begin{theorem}[Semantic globalization requires semantic ledger]
\label{thm:philosophy-semantic-globalization-requires-semantic-ledger}
A natural-language concept used as a global explanation requires a
semantic global ledger.
\end{theorem}
\begin{proof}
Natural-language concepts have ambiguous boundaries, domain dependence,
and unsupported connotations. A global use of such a concept must record
which behaviors fall under it, which do not, which are ambiguous, and
where the concept is not claimed. Those rows are exactly the semantic
global ledger.
\end{proof}

\subsection{Formal and Diagonal Residue}
\label{subsec:philosophy-formal-diagonal-residue}

\begin{definition}[Formal globalizer]
\label{def:philosophy-formal-globalizer}
A formal globalizer uses a checked judgment
$$
\mathcal V\vdash\pi:\varphi
$$
as a claim about the model or behavior family beyond the scope of the
statement $\varphi$.
\end{definition}

\begin{definition}[Verification global ledger]
\label{def:philosophy-verification-global-ledger}
A verification global ledger records checked statements, their scopes,
unverified components, backend, trust boundary, proof dependencies,
translation assumptions, portability boundary, and behaviors not covered
by the verification claim.
\end{definition}

\begin{theorem}[Scoped proof does not imply global explanation]
\label{thm:philosophy-scoped-proof-not-global-explanation}
A formally checked statement on scope $\Omega$ does not imply a global
explanation unless the checked statement and its verification ledger are
global.
\end{theorem}
\begin{proof}
The proof establishes the formal statement that was checked. If the
statement is scoped, its conclusion is scoped. A larger explanation
claim requires additional behavior coverage, transcription ledgers,
verification ledgers, and non-hardenable residue rows.
\end{proof}

\begin{definition}[Diagonal residue]
\label{def:philosophy-diagonal-residue}
For an explanation $E$, diagonal residue is a behavior case, boundary
case, input, intervention, counterexample, out-of-distribution case, or
semantic ambiguity that lies in a not-covered region or ledger gap of
$E$.
\end{definition}

\begin{definition}[Self-exposure test]
\label{def:philosophy-self-exposure-test}
A self-exposure test for $E$ searches for unrecorded source rows,
unrecorded boundaries, false positives, false negatives, semantic
ambiguity, untested interventions, verification gaps, and generalization
residue.
\end{definition}

\begin{theorem}[Non-global ledgers admit diagonal-residue risk]
\label{thm:philosophy-non-global-ledgers-admit-diagonal-residue-risk}
If an explanation lacks a global ledger, then diagonal residue remains a
possible risk.
\end{theorem}
\begin{proof}
Without a global ledger, the required residue of the global claim is not
known to be covered by the recorded ledger rows. Any uncovered required
row may be exposed by an input, behavior case, intervention, or semantic
ambiguity. The absence of the global ledger is therefore a diagonal
residue risk even when a specific witness has not yet been displayed.
\end{proof}

\begin{theorem}[Diagonal residue forces classifier refinement or ledger repair]
\label{thm:philosophy-diagonal-residue-forces-refinement-or-ledger-repair}
When diagonal residue changes the boundary of the supporting classifier,
the explanation must either refine the classifier or repair the ledger.
\end{theorem}
\begin{proof}
If the residue shows that the old classifier wrongly joins or separates
cases, the classifier relation must change. If the classifier is kept,
the residue must be recorded as a boundary, failure case, or not-claimed
row. These are respectively classifier refinement and ledger repair.
\end{proof}

\subsection{Cost and Free Explanation}
\label{subsec:philosophy-cost-free-explanation}

\begin{definition}[Global explanation cost]
\label{def:philosophy-global-explanation-cost}
For a global behavior family, the global explanation cost is the cost of
its NameCert rows, ledger rows, verification rows where used, and seal
rows. When the behavior family is finite this can be written as
$$
\operatorname{Cost}_{\mathrm{global}}(M_\theta)
=
\sum_{B\in\mathcal B(M_\theta)}
\bigl(
\operatorname{NameCertCost}(B)
+
\operatorname{LedgerCost}(B)
+
\operatorname{SealCost}(B)
\bigr).
$$
For nonfinite behavior families, the sum is replaced by the appropriate
cost observer for the declared scope.
\end{definition}

\begin{definition}[Global explanation compression]
\label{def:philosophy-global-explanation-compression}
Global explanation compression uses a smaller family of schemas to cover
many behavior explanations. The compression is certified only when its
schema ledger records coverage, failure cases, uncovered behavior
classes, schema-specific debt, and the bridge from schemas to behaviors.
\end{definition}

\begin{theorem}[Global explanation has nonzero certificate cost]
\label{thm:philosophy-global-explanation-nonzero-cost}
If the covered behavior family contains a nontrivial behavior, then a
global explanation has nonzero certificate cost.
\end{theorem}
\begin{proof}
A nontrivial covered behavior requires at least a supporting classifier,
certificate rows, ledger rows, and a seal. These rows are nonempty
audit objects. Their construction and maintenance cost is therefore not
zero.
\end{proof}

\begin{theorem}[Compressed global explanations require schema ledgers]
\label{thm:philosophy-compressed-global-explanations-require-schema-ledgers}
A schema that compresses many behavior explanations requires a schema
ledger.
\end{theorem}
\begin{proof}
Compression hides behavior-specific boundaries, failure cases, domain
differences, and residue. The schema ledger records which cases the
schema covers, which it does not, and which bridge rows connect the
schema to each behavior explanation. Without those rows the compressed
global claim is unclosed.
\end{proof}

\begin{definition}[Free explanation and free gain]
\label{def:philosophy-free-explanation-free-gain}
A free explanation claims that $E$ explains $B$ without supplying a
classifier, NameCert, ledger, scope, and not-claimed boundary. A free
global explanation makes the same omission for an entire behavior
family. A free interpretability-gain claim asserts positive explanation
information without subtracting certificate cost, ledger cost,
verification cost, and remaining debt.
\end{definition}

\begin{theorem}[No free explanation]
\label{thm:philosophy-no-free-explanation}
A free explanation is not a closed explanation.
\end{theorem}
\begin{proof}
Closed explanation requires a supporting classifier, certificate rows,
ledger rows, scope seal, and not-claimed boundary. A free explanation
omits these rows by definition, so it lacks the closure conditions.
\end{proof}

\begin{theorem}[No free global explanation]
\label{thm:philosophy-no-free-global-explanation}
A free global explanation is not a valid global resolution claim.
\end{theorem}
\begin{proof}
A free global explanation is already missing the rows required for
ordinary closed explanation, and global resolution additionally requires
global behavior coverage and a global ledger. It therefore lacks both
local closure rows and global closure rows.
\end{proof}

\begin{theorem}[No free interpretability gain]
\label{thm:philosophy-no-free-interpretability-gain}
An interpretability-gain claim that omits certificate cost, ledger cost,
verification cost, or remaining debt is not audit-complete.
\end{theorem}
\begin{proof}
Explanation information is benefit minus cost and debt. Omitting any
required cost or debt term overstates the value of the explanation and
prevents an auditor from determining whether the net gain is positive.
\end{proof}

\begin{theorem}[No global black-box eliminator theorem]
\label{thm:philosophy-no-global-black-box-eliminator-theorem}
There is no uncertified global black-box eliminator. Machine-learning
black-box resolution is scoped, classifier-relative, NameCert-bound, and
ledger-complete unless a global behavior family, classifier family,
NameCert family, global ledger, verification ledger where applicable,
and global seal are displayed.
\end{theorem}
\begin{proof}
\autoref{thm:philosophy-scoped-resolution-not-global-resolution} and
\autoref{thm:philosophy-finite-sample-closure-not-open-domain-closure}
separate scoped and finite closure from global closure.
\autoref{thm:philosophy-global-resolution-requires-global-ledger-principle}
gives the global-ledger requirement, and
\autoref{thm:philosophy-no-uncertified-global-black-box-eliminator}
rules out eliminator claims that omit it.
\autoref{thm:philosophy-single-classifier-not-many-behaviors},
\autoref{thm:philosophy-product-classifiers-require-bridge-ledgers},
and
\autoref{thm:philosophy-semantic-globalization-requires-semantic-ledger}
rule out unsupported classifier and semantic globalization.
\autoref{thm:philosophy-scoped-proof-not-global-explanation}
rules out proof-scope overreach.
\autoref{thm:philosophy-non-global-ledgers-admit-diagonal-residue-risk}
and
\autoref{thm:philosophy-diagonal-residue-forces-refinement-or-ledger-repair}
identify the diagonal-residue risk.
\autoref{thm:philosophy-global-explanation-nonzero-cost},
\autoref{thm:philosophy-compressed-global-explanations-require-schema-ledgers},
\autoref{thm:philosophy-no-free-explanation},
\autoref{thm:philosophy-no-free-global-explanation}, and
\autoref{thm:philosophy-no-free-interpretability-gain}
give the cost and audit requirements. Together these statements prove
the theorem.
\end{proof}
